\subsection*{Source-local stress without a separate future-tail input (22 October 2026)}
Let \(y\) denote the source's logarithmic radial coordinate and let the unedited azimuthal profile \(E_0(y)\) be independent of \(\eta\) on and after the A.11 hull of length \(23/10\). The actual two-bump correction has
\(E=E_0g\), \(g=1+\sum_{i=1}^2 c_i(\eta)\beta_i(y)\).
Since the prescribed supports are disjoint and
\(\|\beta_i\|_\infty\le640/3\), put
\(a=(640/3)\max_i\|c_i\|_\infty\) and
\(b=(640/3)\max_i\|\partial_\eta c_i\|_\infty\).
Then \(|g-1|\le a\), \(|g_\eta|\le b\), including on the intervening gap.

Under the \emph{exact same-profile} terminal gauges
\(S(\infty,\eta)=\Pi(\infty,\eta)=0\), source (4.16) reads
\[
N_s=2\eta(hF_1-AF_0)-\tfrac{1-\eta^2}{2}
\partial_\eta(F_1-F_0),\quad A=\tfrac12+h,
\]
with \(F_0=\int_y^\infty E(v)^2\,dv\) and
\(F_1=\int_y^\infty e^{v-y}E(v)^2\,dv\).
The prescribed future slopes give
\(\ell_0=\partial_y\log(\sqrt X E_0)\le0\) and
\(F_{1,0}/E_0(y)^2\le93/20+2h^7/3<5\);
also \(F_{0,0}/E_0(y)^2\le1\).
Crucially \(E_0\) is independent of \(\eta\) on the entire future;
thus \(\partial_\eta(F_1-F_0)\) is supported only between
the current point and the right A.11 support edge, and
\[
|\partial_\eta(F_1-F_0)|
\le2(1+a)b E_0(y)^2
\int_0^{23/10}(1-e^{-t})dt
\le(46/10)(1+a)bE_0(y)^2.
\]
Consequently the complete closed hull satisfies
\begin{equation}\label{eq:A11new-tail}
\sup_{y,\eta}\frac{|N_s|}{E}
\le\frac{E_*}{1-a}
\left[2(1+a)^2(A+5h)+(23/10)(1+a)b\right],
\quad E_*=\sup_{\mathrm{hull}}E_0.
\end{equation}
This retains the angular derivative, and requires no isolated
numerical future-tail \(N_s/E\) input.

The source stage order yields
\(E_*\le e^{47/40}P_*e^{-T_d/2}\lambda^{45}
e^{-13/(2\lambda)}
<8P_*e^{-T_d/2}\lambda^{45}e^{-13/(2\lambda)}\).
After choosing the already source-permitted small parameter
\[
0<\lambda\le\min\{1/100,[100(1+\log P_*)]^{-1}\},
\qquad 0<h\le\lambda/100,
\]
this implies \(E_*<\lambda/1000\). If the authentic A.11 root and
parameter derivative have been independently verified to obey
\(|c_i|\le2\lambda^{28}\),
\(|c_{i,\eta}|\le\lambda^{28}\),
then \(a,b<1/1000\) and
\[
\frac{2(1001/1000)^2(2503/5000)
 +(23/10)(1001/1000)(1/1000)}{999/1000}
=\frac{2513764253}{2497500000}<\frac{11}{10}.
\]
Equation \eqref{eq:A11new-tail} therefore proves
\[
\sup_{\mathrm{entire\ A.11\ hull},|\eta|\le1}|N_s|/E
<11\lambda/10000<\lambda/500<\lambda/2.
\]

\paragraph{Dependency boundary.}
This improves the \emph{conditional} whole-hull stress theorem,
but the authentic moment root and exact same-profile infinity
\(S,\Pi\) gauges remain prerequisites. It does not certify all
exterior stress patches, Appendix-B nonlinear contraction,
the unforced similarity-order startup equations or the physical
PDE. The originating OpenAI whole-space construction has forcing.
The complete proof with definitions, integration and source
references is in the standalone 22 October addendum.